\section{Limitations and Scope}
\label{sec:limitations}

\textbf{Scale of the runs.} These are first, small runs: one calibration attempt per job and option, ten projects for
the replayed allocators and six for the live ones (twice each for \sys{} v2). Replay draws on single attempts, so it
carries none of the models' own variance; the intervals we give are over projects, and with six projects they are wide.

\textbf{The calibration.} Three of the seven validity checks fail on this first table. Knowing each job's hidden
difficulty is worth under one point, so these runs cannot show what adapting during the task is worth, for any
allocator; and the dear model leaves headroom only on long aggregation jobs. A harder calibration is the first item on
the roadmap.

\textbf{Live against replay.} The baselines are replayed with each job's canonical prompt, and under the older rule that
charges a call in full even past the pool. \sys{} v2 ran live, wrote its own messages to its hires, and was held to its
pool at every call. The comparison therefore favors the replayed baselines on both counts.

\textbf{The market and the system.} Prices are list prices on the run date, and the runs have two model tiers, so local
models (and memory as a second currency) are tested in \sys{} but not measured here. Declared difficulty is the agent's
own judgment. One API key is one trust domain, and the hub's token estimate can let a run end a few percent over its
budget (2\% in one of twelve runs).
